\documentclass[10pt,a4paper]{amsart}
\usepackage{amsmath,amssymb,amsthm,graphicx}
% For commutative diagrams we use Xy-pic
\usepackage[all]{xy}

% We use 2cell for 2-commutative diagrams.
\xyoption{2cell}
% Theorem environments.
%
\theoremstyle{plain}
\newtheorem{theorem}[subsection]{Theorem}
\newtheorem{proposition}[subsection]{Proposition}
\newtheorem{lemma}[subsection]{Lemma}

\theoremstyle{definition}
\newtheorem{definition}[subsection]{Definition}
\newtheorem{example}[subsection]{Example}
\newtheorem{exercise}[subsection]{Exercise}
\newtheorem{situation}[subsection]{Situation}

\theoremstyle{remark}
\newtheorem{remark}[subsection]{Remark}
\newtheorem{remarks}[subsection]{Remarks}

\numberwithin{equation}{subsection}

% Macros
%
\def\lim{\mathop{\mathrm{lim}}\nolimits}
\def\colim{\mathop{\mathrm{colim}}\nolimits}
\def\Spec{\mathop{\mathrm{Spec}}}
\def\Hom{\mathop{\mathrm{Hom}}\nolimits}
\def\Ext{\mathop{\mathrm{Ext}}\nolimits}
\def\SheafHom{\mathop{\mathcal{H}\!\mathit{om}}\nolimits}
\def\SheafExt{\mathop{\mathcal{E}\!\mathit{xt}}\nolimits}
\def\Sch{\mathit{Sch}}
\def\Mor{\mathop{\mathrm{Mor}}\nolimits}
\def\Ob{\mathop{\mathrm{Ob}}\nolimits}
\def\Sh{\mathop{\mathit{Sh}}\nolimits}
\def\NL{\mathop{N\!L}\nolimits}
\def\CH{\mathop{\mathrm{CH}}\nolimits}
\def\proetale{{pro\text{-}\acute{e}tale}}
\def\etale{{\acute{e}tale}}
\def\QCoh{\mathit{QCoh}}
\def\Ker{\mathop{\mathrm{Ker}}}
\def\Im{\mathop{\mathrm{Im}}}
\def\Coker{\mathop{\mathrm{Coker}}}
\def\Coim{\mathop{\mathrm{Coim}}}

% Boxtimes
%
\DeclareMathSymbol{\boxtimes}{\mathbin}{AMSa}{"02}

%
% Macros for moduli stacks/spaces
%
\def\QCohstack{\mathcal{QC}\!\mathit{oh}}
\def\Cohstack{\mathcal{C}\!\mathit{oh}}
\def\Spacesstack{\mathcal{S}\!\mathit{paces}}
\def\Quotfunctor{\mathrm{Quot}}
\def\Hilbfunctor{\mathrm{Hilb}}
\def\Curvesstack{\mathcal{C}\!\mathit{urves}}
\def\Polarizedstack{\mathcal{P}\!\mathit{olarized}}
\def\Complexesstack{\mathcal{C}\!\mathit{omplexes}}
% \Pic is the operator that assigns to X its picard group, usage \Pic(X)
% \Picardstack_{X/B} denotes the Picard stack of X over B
% \Picardfunctor_{X/B} denotes the Picard functor of X over B
\def\Pic{\mathop{\mathrm{Pic}}\nolimits}
\def\Picardstack{\mathcal{P}\!\mathit{ic}}
\def\Picardfunctor{\mathrm{Pic}}
\def\Deformationcategory{\mathcal{D}\!\mathit{ef}}
\newenvironment{reference}{}{}
\newenvironment{slogan}{}{}
\newenvironment{history}{}{}
\usepackage{cleveref}
\crefname{originallemma}{original lemma}{original lemmas}
\Crefname{originallemma}{Original lemma}{Original lemmas}
\makeatletter
\newlabel{resolve-lemma-regular-dim-2-projective}{{0C5P}{source}{}{sourcepointer-1}{}}
\newlabel{resolve-lemma-regular-dim-2-projective@cref}{{[originallemma][0][]0C5P}{[1][0][]source}{}{}{}}
\makeatother
\crefname{originallemma}{original lemma}{original lemmas}
\Crefname{originallemma}{Original lemma}{Original lemmas}
\makeatletter
\newlabel{more-morphisms-lemma-projective}{{0B45}{source}{}{sourcepointer-2}{}}
\newlabel{more-morphisms-lemma-projective@cref}{{[originallemma][0][]0B45}{[1][0][]source}{}{}{}}
\makeatother
\crefname{originalsection}{original section}{original sections}
\Crefname{originalsection}{Original section}{Original sections}
\makeatletter
\newlabel{resolve-section-minus-one}{{0C2I}{source}{}{sourcepointer-3}{}}
\newlabel{resolve-section-minus-one@cref}{{[originalsection][0][]0C2I}{[1][0][]source}{}{}{}}
\makeatother
\crefname{originallemma}{original lemma}{original lemmas}
\Crefname{originallemma}{Original lemma}{Original lemmas}
\makeatletter
\newlabel{resolve-lemma-contract-ample}{{0C2M}{source}{}{sourcepointer-4}{}}
\newlabel{resolve-lemma-contract-ample@cref}{{[originallemma][0][]0C2M}{[1][0][]source}{}{}{}}
\makeatother
\title{A regular proper model of a curve is reached from a minimal model by successively contracting exceptional curves of the first kind}
\author{The Stacks Project Authors}
\date{stacks-project:models.tex:local-snapshot-20261012; GNU FDL 1.2 or later}
\begin{document}
\maketitle
\noindent\textit{Editorial note.} The counted unit is the Stacks Project Lemma lemma-pre-exists-minimal-model at source lines 4484-4499 together with its proof at source lines 4501-4516, a straight adjacent pairing in the chapter file. The chapter is the upstream file models.tex; the pool directory carries the longer name stacks-models and is not used as the chapter name anywhere in this card.\par\medskip
\noindent\textit{Editorial note.} The statement opens with the author's slogan environment, which the retained wrapper reduces to a visible no-op; the slogan is the author's own one-line summary of the lemma and is retained verbatim inside the statement span. It is part of the counted unit, not a separate one.\par\medskip
\noindent\textit{Editorial note.} This card has no context span: the statement and proof are self-contained apart from the four pointers below, and the notions they use are recorded in the standard prerequisites.\par\medskip
\noindent\textit{Editorial note.} The retained proof names four results of two other chapters, the projectivity theorem for regular two-dimensional schemes, the existence of an ample invertible sheaf, the section defining exceptional curves of the first kind, and the contraction theorem for such a curve. All four are declared as non-load-bearing navigation pointers with their stable Stacks tags looked up in the shared tag index; none of their statements is reproduced inside the span.\par\medskip
\noindent\textit{Editorial note.} No citation occurs in the retained spans, so no bibliography entry is supplied.\par\medskip
\noindent\textit{Allowed standard prerequisites (editorial).}\begin{itemize}
\item Regular proper models of a curve, exceptional curves of the first kind and their contraction. A model of a smooth projective curve over the fraction field of a discrete valuation ring is a flat proper scheme over the ring with that curve as generic fibre; over a two-dimensional base such a scheme is projective once it is regular, and a projective scheme carries an ample invertible sheaf. An exceptional curve of the first kind is a smooth rational curve of self-intersection minus one inside the surface, and such a curve may be contracted to a regular projective surface, which decreases the number of irreducible components of the special fibre by exactly one. The author's own notions and the three results of the Resolution of Surfaces chapter and the one of More on Morphisms that the proof quotes are the allowed inputs.. The counted lemma is the existence half of the theory of minimal models in this chapter: it produces a minimal model from any regular proper one, and the chapter's later results on uniqueness of minimal models take its conclusion as their starting point.
\end{itemize}\medskip
\ifdefined\hypertarget\hypertarget{sourcepointer-1}{}\fi
\noindent\textit{Source reference.} Original lemma 0C5P: focused-seeds/AGP/sources/stacks-snapshot-20261011/tags.txt, source lines 15269--15269. Allowed input of the Resolution of Surfaces chapter quoted by the retained proof for the projectivity of a regular proper model over a two-dimensional base; the pointer is registered so the reference resolves and the cited statement is not reproduced inside the span. The number is the stable Stacks tag of the target, looked up in the shared tag index.\par\medskip
\ifdefined\hypertarget\hypertarget{sourcepointer-2}{}\fi
\noindent\textit{Source reference.} Original lemma 0B45: focused-seeds/AGP/sources/stacks-snapshot-20261011/tags.txt, source lines 13988--13988. Allowed input of the More on Morphisms chapter quoted by the retained proof for the existence of an ample invertible sheaf on a projective scheme; the pointer is registered so the reference resolves and the cited statement is not reproduced inside the span. The number is the stable Stacks tag of the target, looked up in the shared tag index.\par\medskip
\ifdefined\hypertarget\hypertarget{sourcepointer-3}{}\fi
\noindent\textit{Source reference.} Original section 0C2I: focused-seeds/AGP/sources/stacks-snapshot-20261011/tags.txt, source lines 15158--15158. Allowed input of the Resolution of Surfaces chapter quoted by the retained proof as the place where exceptional curves of the first kind are defined; the pointer is registered so the reference resolves and the cited material is not reproduced inside the span. The number is the stable Stacks tag of the target, looked up in the shared tag index.\par\medskip
\ifdefined\hypertarget\hypertarget{sourcepointer-4}{}\fi
\noindent\textit{Source reference.} Original lemma 0C2M: focused-seeds/AGP/sources/stacks-snapshot-20261011/tags.txt, source lines 15162--15162. Allowed input of the Resolution of Surfaces chapter quoted by the retained proof for the contraction of an exceptional curve by a morphism that keeps the surface regular and projective; the pointer is registered so the reference resolves and the cited statement is not reproduced inside the span. The number is the stable Stacks tag of the target, looked up in the shared tag index.\par\medskip
\begin{lemma}
\label{lemma-pre-exists-minimal-model}
\begin{slogan}
A regular proper model of a curve is obtained by successive blowups
from a minimal model
\end{slogan}
Let $C$ be a smooth projective curve over $K$ with
$H^0(C, \mathcal{O}_C) = K$. If $X$ is a regular proper
model for $C$, then there exists a sequence of morphisms
$$
X = X_m \to X_{m - 1} \to \ldots \to X_1 \to X_0
$$
of proper regular models of $C$, such that each morphism is a
contraction of an exceptional curve of the first kind, and such
that $X_0$ is a minimal model.
\end{lemma}

\begin{proof}
By Resolution of Surfaces, Lemma \ref{resolve-lemma-regular-dim-2-projective}
we see that $X$ is projective over $R$. Hence $X$ has an ample
invertible sheaf by
More on Morphisms, Lemma \ref{more-morphisms-lemma-projective}
(we will use this below).
Let $E \subset X$ be an exceptional curve of the first kind.
See Resolution of Surfaces, Section \ref{resolve-section-minus-one}.
By Resolution of Surfaces, Lemma \ref{resolve-lemma-contract-ample}
we can contract $E$ by a morphism $X \to X'$ such that $X'$ is
regular and is projective over $R$. Clearly, the number of
irreducible components of $X'_k$ is exactly one less than the
number of irreducible components of $X_k$. Thus we can only
perform a finite number of these contractions until we
obtain a minimal model.
\end{proof}

\end{document}
